\section{Περιγραφή κώδικα μικροελεκτή}

Ο κώδικος του esps32 βασίζεται εμπεριέχει 3 βασικά κομμάτια

\begin{enumerate}
    \item Η βασική δομή του κώδικα, που όπως περιγράψαμε και στο κεφαλαίο για το framework του Arduino,
    βασίζεται στις συναρτήσεις setup() και loop().Τα αλλά δυο κομμάτια εμπεριόχονται μέσα σε αυτές τις δύο συναρτήσεις
    \item Ότι αφορά τις αρχικές ρυθμίσεις του αισθητήρα και την απόκτηση πληροφορίας από αυτόν .Στα πειράματα μας που μαζέψαμε δεδομένα και αναλύσαμε,χρησιμοποιήσαμε μόνο τον BME680,αλλά μελετήσαμε και τον αισθητήρα CSS811 και υπάρχει δυνατότητα πρόσθεσης και άλλων αισθητήρων.
    \item Ότι αφορά με την διασύνδεση του μικροελεκτή στον server,μέσω wifi και http request,για την αποστολή δεδομένων στον server.Επίσης,ο esp32 αποκτά γνώση για την ημερομηνία και ώρα που ισχύει στην στιγμήτης ενεργοποίησης.
\end{enumerate}

Η βασική δομή του κώδικα, αναλύεται στα παραπάνω βήματα.Θα κάνουμε μετά τα απαραίτητα σχόλια γι
\begin{enumerate}
    \item  Αρχικοποιούμε τις μεταβλητές και θέτουμε τις τιμές των σταθερών μεταβλητών σε ειδικό φάκελο.
    
    \item Ο esp32 συνδέεται στα διάφορα wifi δικτύων που βρισκόνται στην λίστα των wifi δικτύων που γνωρίζει το esp32 τον κωδικό τους,
     μέχρι να καταφέρει να συνδεθεί με τον server 
    \item Ο esp32 αναθέτει μοναδικό νούμερο id , μεταξύ των κόμβων αισθητήρων, στον ευατό του, με βάση την MAC διεύθυνση τους.
    \item Τρέχουμε την συνάρτηση αρχικοποίησης για καθένα από τους δυνατούς αισθητήρες που έχουμε,
    με την μόνη συνάρτηση αρχικοποίησης που θα λειτουργήσει να ότι είναι ο αισθητήρα που έχει η συσκευή.
    Αν δεν υπάρχει κάποιος γνωστός αισθητήρας,ο μικροελεκτής μπαίνει συνεχή παύση και ο esp32 αναβοσβήνει το built-in led του NodeMCU-32S development board.
    \item Μέσω συνάρτησης των βιβλιοθηκών της esp32,ο esp32 παίρνει την τωρίνη ημερομηνία και Συντονισμένη Παγκόσμια Ώρα (Coordinated Universal Time,UTC), 
    μέσω αίτησης σε δίκτυο πρωτόκολλο δικτυακού χρόνου (Network Time Protocol servers,NTP servers)
    \item Στην συνάρτηση loop(),σε κάθε επανάληψη της συνάρτησης ελέγχουμε εάν έχει παραχθεί από τον αισθητήρα καινούργια τιμή.
    \item Aν αυτό ισχύει,παίρνουμε την καινούργια αυτή μέτρηση και προσθέτουμε στην μέτρηση αυτή το id του μικροελεκτή,
    το όνομα του αισθητήρα που έχει και την τωρίνη χρονική στιγμή σε ημερομηνία και ώρα (datetime). Μετατρούπε την μέτρηση σε JSON document και την προσθέτουμε στο buffer,
    που αποτελεί ένα JSON array. Ανεβάζουμε τον μετρήτη κατά ένα. Για να μην γεμίσει όλη η RAM με buffered αρχεία, 
    βάζουμε πάνω όριο στον αριθμό των μετρήσεων που μπορούν να κρατηθούν στην μνήμη.
   \item Δίνουμε στον μικροελεκτή την δυνατότητα να πραγματοποιήσει κάποια ενέργεια σχετική με το δίκτυο στην συγκεκριμένη επανάληψη.
    Μετά από οποιαδήποτε σχετική ενέργεια που αφορά την διασύνδεση του esp32 μέσω Wifi,παύουμε οποιαδήποτε δυνατότητα εκτέλεσης τέτοιας ενέργειας, 
    μέχρι να επιτευχθεί καινούργια μέτρηση
    \item Εάν συνεχίζει να υπάρχει η σύνδεση μεταξύ του μικροελεκτή και του server,στέλνουμε στον server μύνημα με όλα τα δεδομένα του buffer 
    και αν η αποστολή του μηνύματος ήταν επιτυχής, τότε όλα τα δεδομένα του buffer καθαρίζονται και η μνήμη απελευθερώνεται 
    \item Εάν η αποστολή του μηνύματος ήταν αποτυχημένη, ή για κάποιο λόγο έχει χαθεί η συνδεση με το wifi,
    ο μικροελεκτής ξανακάνει ένα βήμα την φορά ότι έκανε στην setup συναρτήση,αυτή την φορά όμως σε διαφορετικές φάσεις,λόγω του περιορισμού της μίας εκτέλεσης.
\end{enumerate}   

\subsection{Βασικό πρόβλημα υλοποίησης}

Ο λόγος που ο κώδικας υλοποιήθηκε με αυτό τον τρόπο είναι το γεγονός ότι στην BSEC,η έναρξη της μια λειτουργίας αναγκασμού πρέπει να 
πραγματοποιείται στον ρυθμό δειγματοληψίας που έχει οριστεί. Στο κωδικά μας, ο ρυθμός δειγματοληψίας έχει οριστεί να είναι τα 0,333 Hz,
ή πιο απλά 3 δευτερόλεπτα, τα οποία ορίζονται στις αρχικές ρυθμίσεις του αισθητήρα,μαζί με πόσες ψηφιακές (virtual) εξόδους θέλουμε να έχουμε 
από την βιβλιοθήκη της BSEC για επεξεργασία σημάτων.
Δεν υπάρχει η ανάγκη να υπολογίσουμε εμείς πότε θα δειγματολειπτούμε, καθώς μας δίνεται η επόμενη χρονική στιγμή από την συνάρτηση \texttt{bsec\_sensor\_control()},
η οποία ενεργοποιεί την έναρξη μιας λειτουργίας αναγκασμού,δίνωντας τις σχετικές οδηγίες,και επιστρέφει την επόμενη χρονική στιγμή που θα γίνει δειγματοληψία

Σύμφωνα με τον οδηγό ενσωμάτωσης της BSEC, το χρονικό περιθώριο που έχουμε για να εκτελέσουμε την έναρξη μιας νέας λειτουργίας αναγκασμού, αφού 
έχοντας ξεπεράσει σε χρόνο την στιγμή που μας έδωσε η \texttt{bsec\_sensor\_control()},είναι στα 6.25\% της τιμής περιόδου.
Στην δικιά μας περίπτωση, τα 6.25\% των 3 δευτερολέπτων είναι 0.1875 δευτερόλεπτα, οπότε έχουμε όριο τα 3.1875 δευτερόλεπτα ή τα 2,8125 δευτερόλεπτα.
Εάν αυτό το όριο ξεπεραστεί προς τα κάτω ή προς τα πάνω, τότε η \texttt{bsec\_do\_steps()} επιστρέφει τις τιμές σχετικές με τα VOC,
δηλαδή τις μεταβλητές IAQ index,b-VOC και e-CO2,σταθερές χωρίς καμία αλλαγή και η μεταβλητή που υποδικνύει την ακρίβεια των μετρήσεων ίση με το μηδέν.

Χρησιμοποιώντας τον ενσωματικό κώδικα της bsec για την arduino,δεν θα έχουμε ποτέ πρόβλημα να πέσουμε κάτω από το όριο αυτό, καθώς δεν αφήνει τον μικροελεκτή 
να διατάξει λειτουργία αναγκασμού πριν φτάσουμε στο χρόνο που πρέπει να γίνει η δειγματοληψία. Το πρόβλημα προκύπτει στο να ξεπεράσουμε τον χρόνο αυτό.
Σε συνθήκες που δεν εκτελούσαμε καμιά άλλη διεργασία,αυτό δεν θα ήταν πρόβλημα.Αλλά επείδη ο μικροελεκτής πρέπει να στέλνει και τα δεδομένα στον server,
εάν εκτελείται ακόμα διαδικάσία απόστολης δεδομένων ή διασύνδεσης με το δίκτυο wifi στον χρόνο που πρέπει να εκτελέσουμε την  \texttt{bsec\_sensor\_control}.
Σε ιδανικές συνθήκες αυτό δεν θα ήταν πρόβλημα,καθώς μέσα σε 3 δευτερόλεπτα που έχουμε περιθώριο για την επόμενη δειγματοληψία, 
υπάρχει πολύ χρόνος να σταλεί ένα μύνημα http μερικών byte σε τοπικό δίκτυο.

Όμως,υπάρχει σε ρεαλιστικές περιπτώσεις, το δίκτυο wifi μπορεί να είναι ασταθές ή ο server να έχει διάφορων ειδών αδυναμίες να ανταποκριθεί στην λήψη του μηνύματος.
Αυτές οι περίπτωσεις ήταν πολύ υπάρκτα προβλήματα κατά την διάρκεια εκτέλεσης των πειραμάτων,οπότε υπήρχε η ανάγκη επανασύνδεσης ή επανάληψη αποστολής του μυνήματος,
με αποτέλεσμα να πέφτει πάνω στο χρόνο που είχε οριστεί από την \texttt{bsec\_sensor\_control()} για την επόμενη επανακλήση της.

Παρατηρήθηκε εμπειρικά ότι αν διακοπτόταν η σταθερή δειγματοληψία για λίγο,αλλά μετά επέστρεφε στον κανονικά ρυθμό που όριζε η \texttt{bsec\_sensor\_control()},
η bsec μπορούσε εύκολα να γυρίσει σε ακρίβεια λειτουργίας 2 και 3 από μόνη της, παρόλο που κάποιες δειγματοληψίες έχουν τιμή ακρίβειας 0.
Όμως αν η ανακολουθία αυτή συνεχιστεί, τότε ο αισθήτηρας πρέπει να ξαναυποβληθεί στην διαδικασία calibration που έχει περιγραφεί πιο πάνω. 
Οπότε πρέπει να μειώσουμε όσο γίνεται το χάσιμου χρονισμού της δειγματοληψίας


Η arduino είναι μια γλώσσα προγραμματισμού που επιβάλει την σειριακή εκτέλεση,οπότε δεν είναι εύκολη η παύση μιας διεργασίας για την εκτέλεση κάποιας άλλης.
Υπάρχει η δυνατότητα να σταματήσει η τρέχουσα εκτέλεσης στην Arduino μέσω διακοπής (interrupt) από εξωτερικό pin ή εσωτερικό μετρητή
Η bsec και το bme680 δεν έχει την δυνατότητα  να διατάξει διακοπή, παραμόνο όταν η δειγματοληψία έχει οριστεί σαν Ultra low power (ULP), που η περιόδος δειγματοληψίας είναι τα 5 λεπτά
που μπορεί να διατάξει ο χρήστης άμεση δειγματοληψία μέσω εξωτερικού σήματος, όπως είναι ένα κουμπί.

Βεβαίως,υπάρχουν οι εσωτερικές μετρήτες, που είναι πάντα μια καλή επιλογή και θα μπορούσαν να λειτουργήσουν στην παραπάνω περίπτωση
Ωστόσο,οι Ρουτίνες Εξυπηρέτισηςς Διακοπών (Interrupt Service Routines,ISR) είναι συναρτήσεις που συνιστάται να έχουν μικρό χρόνο εκτέλεσης,
συνήθως αλλάζουν την τιμή κάποιας παγκόσμιας (global) μεταβλητής  για να μην παρεμβάλουν με την γενική λειτουργία του μικροελεκτή.
Και επειδή μεγάλο μέρος από τις συναρτήσεις ενσωμάτωσης της BSEC που χρησιμοποιούμε είναι κομμάτι του Arduino ESP32 και όχι της ίδιας της γλώσσας arduino,
και η Arduino ESP32 είναι βιβλιοθήκες που έχουν φτιαχτεί για να λειτουργούν με την γλώσσα της Arduino, αλλά στην πραγματικότητα καλύπτουν διεργασίες που λειτουργούν σε freeRTOS,
κατασκευασμένες με esp-idf,δεν είναι φρόνιμο να λειτουργήσουμε μια ρουτίνα εξυπηρέτησης τόσο μεγάλης, όσο θα ήταν η συνάρτηση ενσωμάτωσης της BSEC 

Για αυτο το λόγο,επιτρέπουμε να γίνει διαδικασία διασύνδεσης, είτε με τό δίκτυο wifi είτε με τον ίδιο τον server, μόνο αμέσως μετά αφότου έχει πραγματοποιηθεί μια δειγματοληψία.
Και μετά αφότου μπαίνουμε στο μπλοκ του κώδικα που διαχειρίζεται το συγκεκριμένο κομμάτι διασύνδεσης,απαγορεύουμε οποιαδήποτε διασύνδεση μέχρι την επόμενη δειγματοληψία.
Αυτό μας δίνει 3 δευτερολέπτα χρόνου,αρκετά εώς τοτέ να πραγματοποιηθεί οποιαδήποτε ενεργεία σχετική με το δικτύο, έχουμε ορίσει χρόνους διακοπής και timeout που δεν ενεμπλέκονται με το κύκλο της δειγματοληψίας. 
Δυστυχώς,παρατηρήθηκε ότι σε περίπτωση που χανόταν ο υπολογιστής στον οποίο έτρεχε ο server έχανε σύνδεση με το υπάρχον δίκτυο wifi ή δεν δούλευε ο server κατά την διάρκεια μίας αίτησης,
το επιλεγμένο timeout δεν τηρόταν,με αποτέλεσμα ο χρόνος ορίου δειγματοληψίας να παραβιαζότανε. 
Ευτυχώς,αυτές οι περιπτώσεις ήταν λίγες και το λογισμικό της επεξεργασίας σήματος μπορούσε να επαναφέρει την ακρίβεια αρκετά γρήγορα.



Αφιερώσαμε τόσο χρόνο για να εξηγήσουμε το πρόβλημα αυτό είναι γιατί είναι ένα πρόβλημα το οποίο κάποιος μπορεί πολύ εύκολα να συναντήσει κοντά του,
που αν δεν έχει πάρει τον κώδικα από την ίδια την BOSCH,ζητώντας να του αποσταλεί στο προσωπικό του email το πακέτο της BSEC, 
δεν θα είχε πρόσβαση στο integration guide που εξηγεί στο FAQ (Frequently Asked Questions) το παραπάνω πρόβλημα \cite{bsecDownload}.
Η bsec μπορεί να κατέβει σαν βιβλιοθήκη και από την PlatformIO και την Arduino IDE,από τα αντίστοιχα library packages,
όπου περιέχει τον ενσωματικό κώδικα της bsec για την arduino και ότι χρειάζεται για να τρέξει η bsec, αλλά δεν έχει το intergation guide.
Πολλά στοιχεία λειτουργίας την BSEC μπορεί να βρεθούν από τον κώδικα που δίνονται στα πακέτα, 
όπως πως πρέπει να γίνεται το calibration των αισθητήρων που περιλαμβάνεται στο API της βιβλιοθήκης της επεξεργασίας σημάτων της BSEC
αλλά η συγκεκριμένη λεπτομέρεια λειτουργίας της bsec δεν γραφέται κάπου ρητά σε σχόλια.

Οπότε για προγραμματιστές που θα χρησιμοποιήσουν την BSEC μαζί με το bme680 το παραπάνω είναι πολύ σημαντικό.
Κιόλας, οι επόμενες εκδόσεις της BSEC είναι για τον μετέπειτα από τον bme680 αισθητήρα της BOSCH της οικογένεις BMExxx,τον BME688.
Δεν υπάρχουν στις διαθίσιμες βιλιοθήκες που μπορεί κάποιος να κατεβάσει από την PlatformIO ή την Arduino IDE, νεότερη έκδοση της BSEC 
από την έκδοση \texttt{1.8.1492} της BSEC Software Library που χρησιμοποιήσαμε.
Οπότε μπορούμε να θεωρήσουμε ότι αυτή η πληροφορία δεν θα υπάρχει σε κάποιο μελλοντικός προγραμματιστή. 
Το καλό είναι ότι τα πακέτα λογισμικό BSEC δεν χρειάζεται χρήματα,λογαριασμό ή κάποια αναμονή για να τα ζητήσεις από το site της BOSCH και η φόρμα δεν ελέγχεται για τα στοιχεία της,πέρα από το email.

Εδώ,θα θέλαμε να σχολιάσουμε ότι σε περίπτωση που δουλεύαμε με την freeRTOS μέσω esp\-idf,θα είχαμε την δυνατότητα να διαχειριστούμε πολλαπλές διεργασίες παράλληλα, 
θα μπορούσαμε να διαχειριστούμε πολύ πιο εύκολα τις διεργασίες παράλληλα,με την δυνατότητα multitasking που έχει η freeRTOS.
Όμως,η βιβλιοθήκη ενσωμάτωσης της BSEC είναι υλοποιημένη μόνω στην Arduino,οπότε θα χρειαζόταν να αναπτύξουμε είτε δικιές μας συναρτήσεις 
για να χρησιμοποιήσουμε την βιβλιοθήκη επεξεργασίας σημάτων της BSEC και την BME68X Sensor API, είτε να υλοποιούσαμε την βιβλιοθήκη ενσωμάτωσης της BSEC σε freeRTOS,
καθώς η κοινότητα για την fre eRTOS και την esp\-idf είναι αρκετά μικρότερη από την Arduino.
Αυτό αντιστοιχεί στα θετικά και αρνητικά της κάθε επιλογής για κάποιον που θα ήθελε να προγραμματίζει σε freeRTOS.

\subsection{Λεπτομέρειες σχετικά με τον κώδικα που αφορά τους αισθητήρες και το buffer}

Στο κομμάτι των αρχικών ρυθμίσεων του αισθητήρα,όπως αναφέρθηκε, συνδέσαμε τον bme680 με τον esp32 μικροελεκτή με to I2C προτόκωλο.
Ο λόγος που επιλέχθηκε αυτό το προτώκολο επικοινωνίας σε σχέση με την SPI,
είναι ότι δεν χρειαζόμαστε την μεγαλύτερη ταχύτητα που προσφέρει o SPI,
χρειάζονται μόλις δύο καλώδια για την διάσυνδεση με τον I2C,σε αντίθεση των 4 που έχει ο SPI και είναι η προεπιλογή στο παράδειγμα που δίνεται
στην βιβλιοθήκη ενσωμάτωσης,οπότε μπορούμε να θεωρήσουμε ότι αποτελεί μια σίγουρη προεπιλογή.

Χρησιμοποιούμε την Flash μνήμη του esp32s module,μέσω της βιβλιοθήκης της Arduino EEPROOM, 
για να μπορέσουμε να διατηρήσουμε την πληροφορία των εσωτερικών μοντέλων της BSEC,που συγκεντρώνει κατα
την διάρκεια της λειτουργίας του BME680,για να μπορεί να βελτιώσει τις εξόδους του BME680 σχετικά με την ποιότητα του αέρα.

Δεσμεύονται τα 156 πρώτα bytes της flash μνήμης,ωστέ να μπορούμε να αποθηκεύσουμε το μέγιστο μέγεθος που 
μπορεί να έχει η κατάσταση των εσωτερικών μοντέλων χωρίς να υπάρχουν προβλήματα μνήμης.
Είναι συγκεκριμένος ο αριθμός των εγγραφών που μπορούμε να κάνουμε σε μία flash μνήμη,μεταξύ 10000 και 100000 κύκλων εγγραφών \cite{esp32EEPROM}.
Οπότε περιορίζουμε τον αριθμό των εγγραφών που πραγματοιποιούμε στην flash μνήμη.
Στα παραδείγματα που προσφέρει η βιβλιοθήκη ενσωμάτωσης της BSEC,
αποθηκεύεται η κατάσταση των εσωτερικών μοντέλων,όταν φτάσει η ακρίβεια των εξόδων της BSEC να είναι στο τρία για πρώτη φορά για την συγκεκριμένη περιόδο λειτουργίας του esp32,
Έπειτα,για όσο παραμένει ενεργός ο esp32,η αποθήκευση γίνεται κάθε έξι ώρες συνεχής λειτουργίας.
Δεν έχουμε κάποια γνώση τι ακριβώς αναπαραστούν οι τιμές που αποθηκεύονται.

Στην αρχική ρύθμιση του αισθητήρα και σε κάθε μετέπειτα δειγματοληψία,ελέγχουμε κάθε φορά αν λειτουργεί όπως πρέπει ο αισθητήρος, και στο κομμάτι της λειτουργίας
που αφορά την Sensor API και την BSEC.
Αν εντοπιστεί κάποιου είδος error,ο esp32 οδηγείται σε ατέρμονη λούπα και το υποδικνύει στον χρήστη αναβοσβήνοντας τον ενσωματωμένο LED του esp32 development board.

Όταν ενημερωνόμαστε από την βιβλιοθήκη ενσωμάτωσης ότι έχουμε καινούργια τιμή δειγματοληψίας από την Api Sensor και έχει επεξεργαστεί από την BSEC,
τότε αποθηκεύουμε όλες τις πραγματικές και ψηφιακές εξόδες,δημιουργημένες από την BSEC,σε ένα \texttt{JsonDocument}.

Η ArduinoJson αποτελεί ανοιχτού κώδικα βιβλιοθήκη \cite{arduinojson},που μας δίνει την δυνατότητα να αποθηκεύσουμε 
όλες τις εξόδους της δειγματοληψίας σε ένα αντικείμενο που έχει μορφή δεδομένων JSON (JavaScript Object Notation).
Ένα json document είναι πολύ εύκολο να αποσταλθεί σε http POST αίτηση και 
να το διαχειριστεί οποιδήποτε λογισμικό το παραλάβει,διότι η json δεν εξαρτάται από κάποια γλώσσα προγραμματισμού
και είναι πολύ ευρεία σε χρήση \cite{json}.

Ακόμα, η ArduinoJson έχει δυνατότητα δυναμικής δέσμευσης μνήμης με την χρήση της \texttt{JsonArray} κλάση,
οπού είναι μία δομή πινάκων \texttt{JsonDocument} αντικειμένων, ενθυλακομένωμη επίσης σε ένα  \texttt{JsonDocument} αντικείμενο
Με αυτόν τον τρόπο υλοποιούμε, και τον buffer για τις μετρήσεις μας, σε περίπτωση που δεν έχουμε συνδέση με τον server για να στείλουμε άμεσα το μύνημα μας,
κρατούμε όσες μετρήσεις έχουμε και τις στέλνουμε όλες μαζί όποτε έχουμε την ευκαιρία.
Για αυτό και αποθηκεύουμε το datetime της στιγμής που πραγματοποιήθηκε η μέτρηση, ωστέ να μπορούμε να διαχωρίσουμε πότε έγινε η κάθε μέτρηση όταν τις παραλάβουμε στον server,
χωρίς παραπάνω δουλειά από τον server να ξεχωρίσει σε ποια στιγμή η κάθε μέτρηση δειγματοληπτήθηκε

Η ArduinoJson έχει το κακό ότιτα \texttt{JsonDocument} χρησιμοποίει πιο πολύ χώρο σε bytes από ότι θα χρησιμοποιούσε μια απλή δομή struct \cite{arduinoJSONram}.
Στην περίπτωση μας,το μέγεθος του \texttt{JsonDocument} σε json που μετρήθηκε με την χρήση του ArduinoJson Assistant
είναι κάπου στα 300 bytes.
Είναι μεγάλο σαν μεγέθος για  ένα ενσωματωμένο συστημά,αλλά επειδή η esp32 έχει 520 KBytes RAM και λίγα από αυτά χρησιμοποιούνται για την λειτουργία του ίδιου του μικροελεκτή,
έχουμε δυνατότητα αποθήκευσης πολλών δεδομένων χωρίς να προκύψει πρόβλημα μνήμης

Επιλέξαμε για μέγιστο μέγεθος buffer να χωράει 300 μετρήσεις, δηλαδή δεκαπέντε λεπτά μετρήσεων.
Ο αριθμός αυτός ελέχθηκε για την μνήμη που καταλαμβάνει.
Λειτουργεί χωρίς πρόβλημα και αφήνει αρκετή μνήμη διαθέσιμη για μελοντικές επεκτάσεις του λογισμικου 
Αν ξεπεραστούν τα αυτές οι 300 μετρήσεις χωρίς να έχουν αποσταλθεί οι μετρήσεις στον server,
αντικαθιστούμε τις νέες μετρήσεις με τις πιο παλιές. 
Έτσι,έχουμε πάντα τα 15 τελευταία λεπτά των μετρήσεων του αισθητήρα bme680 και των ψηφιακών εξόδων της BSEC.

\subsection{Λεπτομέρειες σχετικά με τον κώδικα που αφορά την διασύνδεση με το δίκτυο wifi και τον server}

Το πρώτο κομμάτι της διασύνδεσης είναι η σύνδεση με το δίκτυο wifi,στο οποίο δίνουμε ρητά τα ονόματα των wifi δικτύων,
μαζί με τους κωδικούς τους,
Η προσθήκη επιπλέων στοιχείων για νέα wifi δίκτυα που μπορεί να θελήσει να γίνει από κάποιο μελλοντικό χρήστη της εφαρμογής,
θα χρειαστεί να προσθέσει τους κώδικους στον πηγαίο κώδικα των μικροελεκτών,δεν υπάρχει κάποια δυνατότητα εισαγωγής στοιχείων wifi
από την σελίδα που αλλιλεπιδρά ο χρήστης του συστήματος.

Διασυνδεόμαστε μεταξύ των ρητά δηλομένων wifi που μπορούν να εντοπιστούν από το esp32 μέχρι να καταφέρει να διασυνδεθεί 
σε κάποιο από αυτό τα wifi δικτύα.Ο χρόνος που δίνουμε για την διασύνδεση του wifi είναι 5 δευτερόλεπτα,που είναι περισσότερο 
από τον χρόνο των 3 δευτερολέπτων που μιλήσαμε παραπάνω.
Ευτυχώς,η σύνδεση με το wifi δίκτυο γίνεται ασυγχρόνα και δεν διακόπτει την ροή του esp32,
μπορούμε να ελέγξουμε μετέπειτα στον προβλεπόμενο χρόνο αν συνδέθηκε ή όχι.

Στην διασύνδεση με τον σέρβερ,πρέπει ο esp32 και κάθε συσκευή που βρίσκονται στο  ίδιο τοπικό δίκτυο wifi, που θέλει να συνδεθεί στον server, να ξέρει τα endpoints του server,το port που ακούει ο server και την ακριβή IP της συσκευής που λειτουργεί τον server.
Απο προεπιλογή, ο router διαμοιράζει τις IP διευθύνησεις στις άλλες συσκευές μέσω DHCP server,οπού δεν είναι απαραίτητα ότι θα μείνουν σταθερές.
Αυτό δεν δημιουργεί πρόβλημα στην καθημερική χρήση απλών χρηστών,καθώς το ρούτερ και τα  Σύστημα Ονοματοδοσίας Τομέων 
(Domain Name System,DNS) αναλαμβάνουν την διαχείρηση των IP διευθυνσένω

Μπορούμε να ορίσουμε στατικές IP στις ρυθμίσεις διασύνδεσης στον υπολογιστή που τρέχει τον server.
Δυστυχώς,τα πειράματα γίνανε σε μέρος που δεν υπήρχε παροχή στο ίντερνετ μέσω router.
Υπήρχε πρόσβαση στο ίντερνετ μέσω των δικτύων κινητής τηλεπικοινωνίας και μέσω του φορητού σημείου πρόσβασης,είχαμε τοπικά δίκτυα wifi. 
Και στο φορητό σημείο πρόσβασης συνεχίζει να υπάρχει δυνατότητα να ορίσεις στατική IP στις ρυθμίσεις το wifi στον υπολογιστή.
Ωστόσο,πρακτίκα αποδείχτηκε ότι το φορητό σημείο σύνδεσης δεν κράταει την στατική ip στην συσκευή,οπότε δεν μπορούμε να βασιστούμε σε αυτή την λύση.

Στο παραπάνω πρόβλημα έδωσε λύση με το να χρησιμοποιήσουμε σε υπολογιστή που τρέχει linux,το πακέτο Avahi \cite{avahi},
για την υλοποίηση του multicast (πολλαπλή αποστολή / εκπομπή) DNS,οπού αναφέρεται εώς mDNS \cite{mDNS}. 
Σαν βασική χρήστη, το mdns επιτρέπει κάθε συσκευή να επιλέξει το δικό της όνομα σαν hostname σε τοπικό επίπεδο και αν αυτό το όνομα δεν υπάρχει είδη στη συγκεκριμένη διασύνδεση τοπικού δικτύου,
μπορεί να συνεχίσει να χρησιμοποιεί το όνομα.
Για να το διαχειριστεί ειδικά το όνομα ο mDNS,πρέπει να δηλώθεί με τη \texttt{.local} επέκταση στο τέλος και να το αποστείλει με multicast IP στα υπολοιπές συσκευές του τοπικού δικτύου.

Έχει πολλές από τις υπάρχουσες λειτουργεικότητες από τα DNS,όπως και από το DNS Service Discovery (DNS-SD) \cite{DNSsd},
και ανοίκει σαν τεχνολογία στις τεχνολογίες δικτύου μηδενικής ρύθμισης (Zero-configuration networking), που επιτρέπει σε συσκευές να δημιουργήσουν το δικό τους δίκτυο,
βασιζόμενες στην οικογένεια πρωτοκόλλων TCP/IP και χωρίς ανάγκη ύπαρξης διαχειριστή.
Το όνομα που δώθηκε αυτόματα και διατηρήθηκε ήταν το \texttt{personal-laptop.local}, 
αλλά μπορεί να αλαγχθεί πολύ εύκολα με τις αναλογές ρυθμίσεις στον υπολογιστή που τρέχει τον server και τον πηγαίο κώδικα των esp32

Αφού βρεθήκε το ανάλογο hostname, προσπαθούμε να εγκαταστήσουμε σταθερή σύνδεση http και tcp μεταξύ  server και client,στέλνοντας μήνυμα σε ειδικό ωστέ να προχωρήσουμε στην αποστολή δεδομένων
Αν δεν εντοπιστεί ττο hostname ή δεν πετύχει η αρχική συνδέσης,αποσυνδεόμαστε από το υπάρχον wifi και ψάχνουμε να συνδεθούμε στο επόμενο wifi,δεν σταματάει πότε ο esp32 να ψάχνει για να διασυνδεθεί με τον server. 
Αν φτάσει στο τέλος της λίστας διαθέσιμών συνδέσεων,ξεκινάει ξανά από την αρχή της λίστας.
Το μπλε φως που βγάζει ο ενσωμετομένος LED του NodeMCU development board μας υποδικνύει ότι βρίσκεται σε λειτουργία αναζήτης λειτουργίας διασύνδεσης,υπάρχει έξτρα από το κόκκικο led που δείχνει ότι ο esp32 board είναι ενεργός.